\paragrafos{Ιδιότητες των τετραπλεύρων}{Β.3.4, «Ιδιότητες παραλληλογράμμου – ορθογωνίου – ρόμβου – τετραγώνου – τραπεζίου – ισοσκελούς τραπεζίου»}

\begin{anagke}
Ο Νίκος έχει ένα μεταλλικό πλαίσιο παραθύρου σε σχήμα ορθογωνίου και θέλει να βεβαιωθεί ότι είναι σωστά στερεωμένο, χωρίς γωνιόμετρο. Η Αγγελική του δείχνει ένα κόλπο: «Μέτρα τις δύο διαγώνιες. Αν είναι ίσες, το πλαίσιο είναι σωστό ορθογώνιο!». Ο Βασίλης δοκιμάζει σε ρόμβο από χαρτόνι: οι διαγώνιες του δεν είναι ίσες, αλλά τέμνονται κάθετα και διχοτομούνται.

Η Μυρτώ διπλώνει ένα χάρτινο παραλληλόγραμμο. Δεν βρίσκει άξονα συμμετρίας, αλλά αν το περιστρέψει μισή στροφή γύρω από το σημείο τομής των διαγωνίων, πέφτει ακριβώς πάνω στον εαυτό του. Η Πολυξένη: «Το κέντρο συμμετρίας του είναι εκεί!».

Σε αυτή την παράγραφο μαθαίνουμε τις ιδιότητες κάθε τετραπλεύρου: πλευρές, γωνίες, διαγώνιες, άξονες και κέντρο συμμετρίας.
\end{anagke}

\begin{lexi}{διαγώνιος, διχοτομούνται}{ονόματα}
Η \lexiword{διαγώνιος} ενώνει δύο μη διαδοχικές κορυφές (διά γωνίας). Δύο τμήματα \lexiword{διχοτομούνται} όταν το καθένα περνά από το μέσο του άλλου.
\end{lexi}

\begin{metafrasi}
Οι διαγώνιες περνούν από το μέσο η μία της άλλης. & διχοτομούνται\\
Οι απέναντι πλευρές είναι ίσες. & $\mathrm{ΑΒ}=\mathrm{ΔΓ}$, $\mathrm{ΑΔ}=\mathrm{ΒΓ}$\\
Οι απέναντι γωνίες είναι ίσες. & $\hat{\mathrm{Α}}=\hat{\Gamma}$, $\hat{\mathrm{Β}}=\hat{\mathrm{Δ}}$\\
Οι διαγώνιες σχηματίζουν ορθή γωνία. & κάθετες: $\mathrm{ΑΓ}\perp\mathrm{ΒΔ}$\\
\end{metafrasi}

\tmima{Η θεωρία}

\begin{kanonas}{παραλληλόγραμμο}
Σε κάθε παραλληλόγραμμο:
\begin{itemize}[leftmargin=1.2em,itemsep=0.3ex,topsep=2pt]
\item το σημείο τομής των διαγωνίων είναι \textbf{κέντρο συμμετρίας},
\item οι διαγώνιες \textbf{διχοτομούνται},
\item οι απέναντι πλευρές είναι \textbf{ίσες},
\item οι απέναντι γωνίες είναι \textbf{ίσες}.
\end{itemize}
\end{kanonas}

\begin{kanonas}{ορθογώνιο}
Εκτός από τις ιδιότητες του παραλληλογράμμου: οι \textbf{διαγώνιες είναι ίσες}, και οι μεσοκάθετοι των πλευρών του είναι \textbf{άξονες συμμετρίας} (δύο άξονες).
\end{kanonas}

\begin{kanonas}{ρόμβος}
Εκτός από τις ιδιότητες του παραλληλογράμμου: οι ευθείες των διαγωνίων είναι \textbf{άξονες συμμετρίας}, οι διαγώνιες είναι \textbf{κάθετες} και \textbf{διχοτόμοι των γωνιών} του.
\end{kanonas}

\begin{kanonas}{τετράγωνο}
Έχει όλες τις ιδιότητες του ορθογωνίου και του ρόμβου: οι διαγώνιες είναι ίσες, κάθετες, διχοτομούνται, διχοτομούν τις γωνίες ($45^\circ$ η καθεμία). Έχει \textbf{$4$ άξονες συμμετρίας} (δύο διαγώνιες και δύο μεσοκάθετοι των πλευρών) και κέντρο συμμετρίας.
\end{kanonas}

\begin{kanonas}{ισοσκελές τραπέζιο}
Η ευθεία που περνά από τα μέσα των βάσεων είναι \textbf{άξονας συμμετρίας} και μεσοκάθετος των βάσεων. Οι \textbf{προσκείμενες στη βάση γωνίες} είναι ίσες.
\end{kanonas}

\begin{center}\small
\begin{tabular}{lcccc}
\toprule
Ιδιότητα & Παραλληλόγραμμο & Ορθογώνιο & Ρόμβος & Τετράγωνο\\
\midrule
Απέναντι πλευρές ίσες & ναι & ναι & ναι & ναι\\
Απέναντι γωνίες ίσες & ναι & ναι & ναι & ναι\\
Διαγώνιες διχοτομούνται & ναι & ναι & ναι & ναι\\
Διαγώνιες ίσες & όχι & ναι & όχι & ναι\\
Διαγώνιες κάθετες & όχι & όχι & ναι & ναι\\
Άξονες συμμετρίας & $0$ & $2$ & $2$ & $4$\\
Κέντρο συμμετρίας & ναι & ναι & ναι & ναι\\
\bottomrule
\end{tabular}
\end{center}

\begin{prosoxi}
\begin{itemize}[leftmargin=1.2em,itemsep=0.3ex,topsep=2pt]
\item Σε ένα (πλάγιο) παραλληλόγραμμο οι διαγώνιες δεν είναι ούτε ίσες ούτε κάθετες, και δεν υπάρχει άξονας συμμετρίας.
\item Ό,τι ισχύει για το παραλληλόγραμμο ισχύει και για το ορθογώνιο, τον ρόμβο και το τετράγωνο, γιατί είναι παραλληλόγραμμα.
\item Οι διαγώνιες του ισοσκελούς τραπεζίου είναι ίσες, αλλά δεν διχοτομούνται.
\end{itemize}
\end{prosoxi}

\tmima{Λυμένα παραδείγματα}

\paradeigma{Το ορθογώνιο παράθυρο}
\tixero Ορθογώνιο παράθυρο με διαστάσεις $90$ cm και $120$ cm.\par
\tithelo Η διαγώνιος και πόσο απέχει το κέντρο συμμετρίας από μια κορυφή.\par
\possindeo Από το πυθαγόρειο: $\mathrm{ΑΓ}=\sqrt{90^2+120^2}=\sqrt{8100+14400}=150$ cm. Οι διαγώνιες διχοτομούνται και είναι ίσες, άρα το κέντρο απέχει $150:2=75$ cm από κάθε κορυφή.\par
\apantisi $150$ cm και $75$ cm.\par
\elegxos Οι τέσσερις αποστάσεις από το κέντρο είναι ίσες.

\paradeigma{Οι διαγώνιες του ρόμβου}
\tixero Ρόμβος με διαγώνιες $10$ και $24$ cm.\par
\tithelo Η πλευρά του.\par
\possindeo Οι διαγώνιες διχοτομούνται και είναι κάθετες: σχηματίζουν ορθογώνιο τρίγωνο με κάθετες $5$ και $12$. Η πλευρά είναι $\sqrt{5^2+12^2}=13$ cm.\par
\apantisi $13$ cm (περίμετρος $52$ cm).\par
\elegxos $5^2+12^2=25+144=169=13^2$.

\paradeigma{Οι γωνίες του παραλληλογράμμου}
\tixero Παραλληλόγραμμο με $\hat{\mathrm{Α}}=65^\circ$.\par
\tithelo Οι άλλες γωνίες.\par
\possindeo Οι απέναντι είναι ίσες: $\hat{\Gamma}=65^\circ$. Οι διαδοχικές είναι παραπληρωματικές (εντός επί τα αυτά): $\hat{\mathrm{Β}}=180^\circ-65^\circ=115^\circ=\hat{\mathrm{Δ}}$.\par
\apantisi $65^\circ,115^\circ,65^\circ,115^\circ$.\par
\elegxos Το άθροισμα είναι $360^\circ$.

\paradeigma{Το ισοσκελές τραπέζιο}
\tixero Ισοσκελές τραπέζιο με μια γωνία στη μεγάλη βάση $70^\circ$.\par
\tithelo Οι άλλες γωνίες.\par
\possindeo Οι προσκείμενες στη βάση γωνίες είναι ίσες: δύο γωνίες $70^\circ$. Οι γωνίες στη μικρή βάση είναι παραπληρωματικές των πλαγίων: $180^\circ-70^\circ=110^\circ$ η καθεμία.\par
\apantisi $70^\circ,70^\circ,110^\circ,110^\circ$.\par
\elegxos $70+70+110+110=360$.

\tmima{Ασκήσεις}

\zoni{Το κατάλαβα}{πλευρές, γωνίες, διαγώνιοι}

\ex Παραλληλόγραμμο έχει πλευρές $8$ cm και $5$ cm. Βρες τις άλλες δύο και την περίμετρο.

\ex Παραλληλόγραμμο έχει γωνία $50^\circ$. Βρες τις άλλες τρεις γωνίες.

\ex Οι διαγώνιες ορθογωνίου είναι $13$ cm. Πόσο απέχει το κέντρο από κάθε κορυφή;

\ex Ρόμβος έχει διαγώνιες $12$ και $16$ cm. Βρες την πλευρά και την περίμετρο.

\ex Πόσους άξονες συμμετρίας έχει το τετράγωνο; Σχεδίασέ τους.

\ex Σχεδίασε ορθογώνιο ΑΒΓΔ με διάμετρο μια διαγώνιο και δες ότι ο κύκλος περνά από όλες τις κορυφές.

\ex Σωστό ή λάθος; α) Οι διαγώνιες του ορθογωνίου είναι ίσες. \quad β) Οι διαγώνιες του ρόμβου είναι ίσες. \quad γ) Οι διαγώνιες του παραλληλογράμμου διχοτομούνται. \quad δ) Το ορθογώνιο έχει $4$ άξονες συμμετρίας.

\zoni{Το χρησιμοποιώ}{κατασκευές και μετρήσεις}

\ex \textbf{Το πλαίσιο.} Ορθογώνιο πλαίσιο είναι $60$ cm $\times$ $80$ cm. Πόση πρέπει να είναι η διαγώνιος για να είναι σωστό ορθογώνιο;

\ex \textbf{Ο χαρταετός.} Ένας χαρταετός σε σχήμα ρόμβου έχει διαγώνιες $60$ και $80$ cm. Πόση είναι η πλευρά του και πόση η περίμετρός του;

\ex \textbf{Το τραπεζάκι.} Ισοσκελές τραπέζιο έχει γωνία $65^\circ$ στη μεγάλη βάση. Βρες τις άλλες γωνίες του.

\ex \textbf{Το οικόπεδο.} Ένα παραλληλόγραμμο οικόπεδο έχει γωνία $72^\circ$. Βρες τις υπόλοιπες γωνίες.

\ex \textbf{Ο γνώμονας.} Ένας τεχνίτης ελέγχει ότι ένα τραπέζι είναι ορθογώνιο μετρώντας τις διαγώνιες. Αρκεί αυτό αν το τραπέζι είναι σίγουρα παραλληλόγραμμο;

\zoni{Το σκέφτομαι}{γιατί συμβαίνει}

\ex Εξήγησε γιατί οι απέναντι πλευρές παραλληλογράμμου είναι ίσες. (Βοήθεια: δες το κέντρο συμμετρίας.)

\ex Σε ποια παραλληλόγραμμα οι διαγώνιες είναι διχοτόμοι των γωνιών; Εξήγησε.

\ex Βρες τους άξονες συμμετρίας του ισοσκελούς τραπεζίου. Εξήγησε γιατί οι προσκείμενες γωνίες βάσης είναι ίσες.

\ex Ποιο σχήμα προκύπτει αν ενώσουμε διαδοχικά τα μέσα των πλευρών ενός ορθογωνίου; (Ένα πείραμα με χαρτί.)

\zoni{Για το σπίτι}{πείραμα με χαρτί}

\ex Κόψε από χαρτί ένα παραλληλόγραμμο και τσάκισέ το στις διαγώνιες. Τι διαπιστώνεις;

\ex Μέτρησε τις διαγώνιες μιας πόρτας ή ενός παραθύρου και έλεγξε αν είναι ίσες.
